\section{The remaining coefficient-five rows through the cyclotomic class eigenspace}
\label{sec:coefficient-five-cyclotomic-eigenspace}

Put $K=\mathbb Q(\omega)$ and $R=\mathbb Q(\zeta_7+\zeta_7^{-1})$.  For
$D\in\{21,6,14\}$, put
\[
 F_D=R(\sqrt D),\qquad L_D=F_D(\mu_5),\qquad
 \Delta_D=\operatorname{Gal}(L_D/F_D).
 \tag{C5E.1}
\]
Write $\omega_5:\Delta_D\to\mathbb F_5^\times$ for its
action on $\mu_5$.  A superscript $\omega_5$ below means
this natural eigenspace, rather than its inverse.  On a finite
five-primary group use the Teichmuller lift of $\omega_5$ to
$\mathbb Z_5^\times$; on an $\mathbb F_5$-vector space use
$\omega_5$ itself.

For the original coefficient and character convention of AC-G199,
let $O$ be any original ring of integers of a finite extension
$E/\mathbb Q_5$ containing the values of $\xi_7$, where
$\xi_7(3)=\zeta_6$.  Set $D_O=E/O$ and retain the full strict
kernel, including the full localization at five, every other
finite inertia restriction, and the real-place condition.
For $4\mid e$, the nine pairs missing from (F5.1) are
\[
 \mathcal A_e^{\rm rem}
   =\{(1,j):j=1,3,5\}\ \sqcup\
                  \{(\chi_{-8},j):0\leq j<6\}.
 \tag{C5E.2}
\]
Let $\epsilon(\delta)\in\{0,1\}$ satisfy
$\delta(-1)=(-1)^{\epsilon(\delta)}$.  For such a pair set
$b\equiv\epsilon(\delta)+j-e\pmod2$,
$b\in\{0,1\}$, and use the same original rational characters
\[
 \psi^{\rm nc}=\delta\xi_7^j\chi_{-3}^b,
 \qquad a^{\rm nc}=\psi^{\rm nc}c_5^{-e},
 \qquad a^{\rm cr}=a^{\rm nc}\chi_{-3}.
 \tag{C5E.3}
\]

\begin{theorem}[AC-G205: the remaining nine coefficient-five rows]
\label{AC-G205}
For every $D$ in (C5E.1), the maximal pro-five Galois extension
of $F_D$ unramified at all finite places away from five has group
$\mathbb Z_5$.  Each of the two local groups at five maps onto
this group separately.  Its required class input is the exact
eigenspace assertion
\[
               \bigl(\operatorname{Cl}(L_D)\{5\}\bigr)^{\omega_5}=0.
 \tag{C5E.4}
\]
The braces denote the five-primary subgroup.  No assertion of
vanishing of the whole class group, or of its whole odd part,
is included in (C5E.4).

On each original $O$, at every positive weight $e$ divisible
by four and every pair in (C5E.2),
\[
 \begin{aligned}
               \mathcal S_5(\mathbb Q,D_O(a^{\rm nc}))&=0,\\
 \mathcal S_5(K,D_O(\delta\xi_7^j c_5^{-e}))
          &\simeq\mathcal S_5(\mathbb Q,D_O(a^{\rm cr})).
 \end{aligned}
 \tag{C5E.5}
\]
These are the unchanged full strict groups.  The critical group
on the right is retained and is covered by the one aggregate
of AC-G191.

For the actual host, $h\geq2$, and $T\geq1$, the new internally
distinct set is
\[
 \begin{split}
 Y^{\rm rem}_5(T)=\{(5,(\delta\xi_7^j)|_{G_K},e):\;&1\leq e<T,
               \ 4\mid e,\ (\delta,j)\in\mathcal A_e^{\rm rem}\},\\
 |Y^{\rm rem}_5(T)|={}&9\left\lfloor\frac{T-1}{4}\right\rfloor.
 \end{split}
 \tag{C5E.6}
\]
It is disjoint from the G199 family.  Their union contains all
twenty-four pairs $\delta\in\{1,\chi_{-4},\chi_8,\chi_{-8}\}$,
$0\leq j<6$, at every positive weight.  For every retained
set $J$ of original triples,
\[
                 \mathcal R_J
                    =\mathcal R_{J\setminus Y^{\rm rem}_5(T)}.
 \tag{C5E.7}
\]
This deletes zero noncritical summands.  It does not split or
annihilate individual positive derived indices.
\end{theorem}

\begin{proof}
\proofstep{1. Exact real and CM fields, and the two local primes.}
The real cyclic cubic field $R$ has defining polynomial
$X^3+X^2-2X-1$.  Its intersection with each quadratic
$\mathbb Q(\sqrt D)$ is $\mathbb Q$, by the coprime degrees.
Thus $F_D$ is totally real of degree six, with group $C_3\times C_2$.
It is unramified at five, since its cyclotomic moduli are
$21$, $168$, and $56$, respectively.  Its intersection with
$\mathbb Q(\mu_5)$ is trivial: every nontrivial subfield of
the latter is ramified at five.  Consequently
\[
 \begin{gathered}
 [L_D:\mathbb Q]=24,\qquad \Delta_D=C_4,\\
 \operatorname{Gal}(L_D/\mathbb Q)=C_3\times C_2\times C_4.
 \end{gathered}
 \tag{C5E.8}
\]
The maximal real subfield of $L_D$ has degree twelve; it is
not being identified with $F_D$.

In $R$, the Frobenius of five has order three in
$(\mathbb Z/7)^\times/\{1,-1\}$, since $5^3\equiv-1\pmod7$.
The three radicands satisfy $21\equiv6\equiv1$ and
$14\equiv4\pmod5$, with distinct square roots in each case.
Thus $F_D$ has two primes above five, and each completion is
the unramified cubic extension of $\mathbb Q_5$.
The polynomial $\Phi_5(1+X)$ is Eisenstein at five, so the
local cyclotomic extension has totally ramified degree four.
It is disjoint from the unramified cubic completion.
Each prime of $F_D$ therefore has a single extension to $L_D$,
with local pair $(e_5,f_5)=(4,3)$.
Complex conjugation $c\in\Delta_D$ fixes both primes individually.

\proofstep{2. The eighteen required natural class characters.}
Let $k_0=\mathbb F_5[z]/(z^2+z+1)$, and identify the reduction
of $\xi_7^2(3)$ with $z$.  Thus $z$ has order three.
Write $d_D=21,24,56$ for the primitive discriminants of
$\mathbb Q(\sqrt D)$ in the displayed order.  Let $\tau_5$
be the Teichmuller character modulo five, with reduction
$\overline\tau_5(n)=n\pmod5$.  The characters of (C5E.8)
whose restriction to $\Delta_D$ is $\omega_5$ are precisely
\[
        \chi_{D,a,b}=\xi_7^{2a}\chi_{d_D}^{b}\tau_5,
                   \qquad a=0,1,2,\quad b=0,1.
 \tag{C5E.9}
\]
They are odd.  All their values reduce into the fixed field
$k_0$, which is only used to check these finite class components.
This does not replace the coefficient $O$ in (C5E.3).

The exact primitive conductors $f=f_{D,a,b}$ are
\[
 \begin{array}{c|cc}
       &a=0&a=1,2\\\hline
 b=0&5&35\\
 D=21,\ b=1&105&105\\
 D=6,\ b=1&120&840\\
 D=14,\ b=1&280&280
 \end{array}
 \tag{C5E.10}
\]
For $D=21$, the quadratic factor is $\chi_{-3}\chi_{-7}$,
whose seven exponent is three; addition of $2a$ never cancels it.
For $D=6$ the primitive factor $\chi_{24}$ has conductor $24$,
disjoint from seven and five.  For $D=14$ the factor is
$\chi_{-8}\chi_{-7}$, again with a noncancelling odd seven
exponent and primitive conductor eight at two.
These observations prove (C5E.10), including $a=0$.
Every $\varphi(f)$ in the table is prime to five.

\proofstep{3. Integral Stickelberger witnesses and exact norm descent.}
For $Q_f=\mathbb Q(\mu_f)$ and its natural automorphisms
$\sigma_n$, retain the convention
\[
 \begin{aligned}
 \theta_f&=\frac1f\sum_{\substack{1\leq n<f\\(n,f)=1}}
                  n\sigma_n^{-1},\\
 B_{f,c}&=(c-\sigma_c)\theta_f
     =\sum_{\substack{1\leq n<f\\(n,f)=1}}
                  \left\lfloor\frac{cn}{f}\right\rfloor\sigma_n^{-1}.
 \end{aligned}
 \tag{C5E.11}
\]
Here $(c,f)=1$.  The last equality follows by taking the
least positive residue of $cn$ modulo $f$; in particular
the operator is integral before any reduction modulo five.
Sharifi's complete Lemma 4.5.2 and Theorem 4.5.6 show that
it annihilates $\operatorname{Cl}(Q_f)$.

On a natural $\chi=\chi_{D,a,b}$ eigenvector, its scalar is
\[
       S_{D,a,b}(c)=
       \sum_{\substack{1\leq n<f\\(n,f)=1}}
           \left\lfloor\frac{cn}{f}\right\rfloor
                              \overline\chi(n)^{-1}\ \in k_0.
 \tag{C5E.12}
\]
The following exact witnesses are all nonzero; entries are
$(c,S(c))$.  The first row is common to all three $D$.
\[
 \begin{array}{c|ccc}
      &a=0&a=1&a=2\\\hline
 b=0&(2,1)&(2,4+4z)&(2,z)\\
 D=21,\ b=1&(2,3)&(2,2+2z)&(2,3z)\\
 D=6,\ b=1&(7,2)&(11,4+4z)&(11,z)\\
 D=14,\ b=1&(3,3)&(3,2+3z)&(3,4+2z)
 \end{array}
 \tag{C5E.13}
\]
They are finite integer-floor sums with $f\leq840$.
The retained exact checker evaluates (C5E.12) in
$k_0$ using $z^2=-z-1$, checks each primitive conductor,
and records the complete sums and every printed witness.
No floating-point value or uncompleted class computation is used.

For clarity, use the primitive character field $E_\chi$,
the fixed field of the kernel of (C5E.9).
It is contained in both $L_D$ and $Q_f$.
The extension degrees $[L_D:E_\chi]$ and $[Q_f:E_\chi]$
are prime to five.  For either extension $M/E_\chi$, extension
of ideals $j$ and norm $N$ satisfy $N\circ j=[M:E_\chi]$
on the smaller class group and
$j\circ N=\sum_{g\in\operatorname{Gal}(M/E_\chi)}g$ on the larger one.
Thus $j$ injects the five-primary class groups; averaging by
the invertible degree identifies its image with the invariants.
These maps are equivariant.
Consequently a $\chi$ component of
$\operatorname{Cl}(L_D)[5]\otimes_{\mathbb F_5}k_0$ descends
to $E_\chi$ and injects into the corresponding component over
$Q_f$.  The integral annihilator acts there by the nonzero scalar
in (C5E.13).  That component is zero.
This proves vanishing for every character in (C5E.9).
It avoids inserting imprimitive Euler factors into the operator.

Because $24$ is invertible modulo five, the character
decomposition over $k_0$ is exact.  Faithful scalar extension
therefore gives
\[
                   \operatorname{Cl}(L_D)[5]^{\omega_5}=0.
 \tag{C5E.14}
\]
The idempotent for $\omega_5$ on $\Delta_D$ also exists over
$\mathbb Z_5$, using the Teichmuller values and denominator four.
Its image in the finite five-primary class group has no
five-torsion by (C5E.14), hence is zero.  This proves (C5E.4).
The $S$-class group, for $S$ consisting of the two five-primes
and infinity, is an equivariant quotient, so its corresponding
eigenspace is zero as well.  No claim about the inverse
cyclotomic eigenspace has been made.

\proofstep{4. Kummer theory selects the cyclotomic eigenspace.}
Write $U=\mathcal O_{L_D,S}^{\times}$.
For any $u\in U$, the quotient $u/c(u)$ has zero valuation
at every finite prime: outside $S$ both are units, and at
the two primes in $S$ conjugation fixes each prime.
At every complex embedding its absolute value is one.
It is therefore a root of unity.  One direct proof is that
all powers are algebraic integers of bounded degree whose
conjugates have absolute value one; their monic minimal
polynomials have bounded integer coefficients, hence form
a finite set.  Two powers coincide.

In the additive notation of $U/U^5$,
$(1-c)[u]$ is thus in the image of the cyclic group of roots
of unity.  Since two is invertible modulo five, every
minus vector is half such an image.  Consequently
\[
                       \dim_{\mathbb F_5}(U/U^5)^-\leq1.
 \tag{C5E.15}
\]
Use the equivariant $S$-Kummer sequence
\[
 0\longrightarrow U/U^5\longrightarrow H^1(G_{L_D,S},\mu_5)
                    \longrightarrow\operatorname{Cl}(L_D,S)[5]
                    \longrightarrow0.
 \tag{C5E.16}
\]
This is the exact original sequence of Sharifi Proposition 1.4.5.
The common maximal-$S$ extension and prime-to-five restriction
give
\[
 H^1(G_{F_D,S},\mathbb F_5)
             \simeq H^1(G_{L_D,S},\mathbb F_5)^{\Delta_D}.
 \tag{C5E.17}
\]
Choose a primitive fifth root to identify $\mathbb F_5$ with
$\mu_5$ over $L_D$.  For $g\in\Delta_D$, the action on a
degree-one cochain is
$(g f)(\sigma)=g\bigl(f(g^{-1}\sigma g)\bigr)$.
The identification consequently multiplies the constant-coefficient
action by $\omega_5(g)$.  Thus the right side of (C5E.17)
identifies with the $\omega_5$, rather than $\omega_5^{-1}$,
eigenspace of $H^1(G_{L_D,S},\mu_5)$.
The action on the class term in (C5E.16) is the natural class action.
Taking this eigenspace is exact, and (C5E.4) makes the class
term zero.  Since $\omega_5(c)=-1$, (C5E.15) proves
\[
                    \dim_{\mathbb F_5}H^1(G_{F_D,S},\mathbb F_5)\leq1.
 \tag{C5E.18}
\]

\proofstep{5. The global pro-five group and each local image.}
The cyclotomic $\mathbb Z_5$ extension of $\mathbb Q$ is
unramified at every finite place away from five.  Its
intersection with $F_D$ is trivial, because $[F_D:\mathbb Q]=6$
is prime to five.  Thus the maximal pro-five group $P_D$
in the theorem surjects onto $\mathbb Z_5$.
Equation (C5E.18) is therefore an equality.

Every nontrivial finite five-group quotient of $P_D$ has
degree-one $\mathbb F_5$ cohomology of dimension at most one.
Its maximal subgroups are normal of index five and are kernels
of nonzero such homomorphisms.  There is exactly one maximal
subgroup, so any element outside it generates the quotient.
Choose in $P_D$ an element nonzero under its cyclotomic
character modulo five.  In every nontrivial finite quotient
the pullback of any nonzero homomorphism to $\mathbb F_5$
is a nonzero multiple of this unique global character.
The chosen element therefore generates every finite quotient.
It follows that $P_D$ is procyclic and infinite, hence isomorphic
to $\mathbb Z_5$.  Its displayed surjection to $\mathbb Z_5$
is an isomorphism, since a nonzero closed kernel in
$\mathbb Z_5$ would give a finite quotient.

For either $v\mid5$, the degree $[F_{D,v}:\mathbb Q_5]=3$
is prime to every cyclotomic five-power degree.  The local
image therefore surjects onto the cyclotomic $\mathbb Z_5$
quotient, and hence onto $P_D$ itself.  This holds for each
place individually, not just for their combined images.

\proofstep{6. One local witness on the unchanged original coefficient.}
More generally let $a:G_{F_D}\to1+\mathfrak m_O$ be continuous
and unramified away five.  A full strict cocycle
$z:G_{F_D}\to(E/O)(a)$ has finite image.
At one of the five-primes, full local zero supplies a single
coboundary witness $x\in E/O$.  Choose a power of the original
uniformizer killing both $x$ and all values of $z$.
The map $g\mapsto(z(g),a(g))$ into the corresponding finite
semidirect product has five-power image: the additive torsion
module is a five-group, and a finite quotient of $1+\mathfrak m_O$
is a five-group.  At every other finite place strict inertia
zero is literal zero of $z$, since $a$ is unramified there.
Thus this finite image factors through $P_D$.
Its entire global image comes from the chosen local group,
by Step 5.  Each global element consequently has the same
image as a local element, including its scalar action on $x$.
The local coboundary formula with this same $x$ is global.
We have proved original-coefficient strict zero over $F_D$.
The real local condition is automatic on five-primary modules.

For (C5E.2)--(C5E.3), write $a=j\bmod2$.  The exact finite
character decompositions are
\[
 \begin{array}{c|c|c}
 (\delta,j)&\psi^{\rm nc}&D\\\hline
 (1,j),\ j\text{ odd}&\xi_7^{j-3}\chi_{21}&21\\
 (\chi_{-8},j),\ j\text{ even}&\xi_7^{j}\chi_{24}&6\\
 (\chi_{-8},j),\ j\text{ odd}&\xi_7^{j-3}\chi_{56}&14
 \end{array}
 \tag{C5E.19}
\]
Every displayed power of $\xi_7$ is even and is killed by $R$.
The quadratic factors are killed by the corresponding $\sqrt D$.
Thus $F_D$ kills the whole finite character in (C5E.19).
Because $4\mid e$, the restricted cyclotomic coefficient is
near one.  It is unramified at every finite place away from five.
The preceding strict-zero result applies on the original $O$.
Restriction followed by corestriction from $F_D/\mathbb Q$
is multiplication by six, a unit on the five-primary cohomology.
It preserves every full local and inertia condition.
This proves the first assertion in (C5E.5).

\proofstep{7. Original rational descent, actual labels and retained sum.}
The two rational extensions in (C5E.3) restrict to the same
character over $K$.  The two-coset induced original lattice
splits into these two characters with determinant a unit times
two, which is invertible in $O$.  Full Shapiro therefore gives
the exact decomposition of strict groups used in G199 Step 7.
At five, the single completion of $K$ is unramified quadratic;
at three, inertia has the full ramified quadratic quotient;
elsewhere the local double-coset decomposition retains every
inertia condition.  The first rational summand has just been
proved zero, giving the second assertion in (C5E.5).
Only its critical partner remains.

All nine labels in (C5E.2) already belong to the actual host
and prescribed coefficient completion of (FC.1)--(FC.2).
Their finite orders divide $3\cdot2^h$; no conductor-five
character is inserted into a finite host label.
Their rational ratios have conductor dividing $56$, so
distinct pairs remain distinct on $G_K$, whose nontrivial
quadratic kernel character is ramified at three.
The count in (C5E.6) follows, and the literal complementary
lists in (F5.1) prove disjointness from G199.
Every new original-$K$ term is precisely its critical rational
partner, which belongs to the one nonnegative aggregate (CP.2).
For any retained $J$, its noncritical rational terms are zero
at all new pairs, proving (C5E.7) with no extra comparison cost.
Full-union index comparison and same-window transport remain
the separate ensuing step; the all-prime remainder is retained.
\end{proof}

\paragraph{Finite-cover dictionary and exact input scope.}
For every original uniformizer power the finite coefficient
in (C5E.3) is a locally constant etale sheaf on
$\operatorname{Spec}\mathbb Z[1/(2\cdot3\cdot5\cdot7)]$,
with the stated completion, inertia and real-place restrictions.
The auxiliary covers $F_D$ and $L_D$ trivialize the stated
residual actions and make the Kummer coefficient character explicit.
This retains the SGA~1, Expose V, sections~4 and~7 finite-cover
dictionary at the preceding (FC.2) audit scope; no new archival
reading or motivic realization is asserted.
Root proposed the two-prime CM route and explored full class
computations.  The present author supplies the targeted cyclotomic
eigenspace replacement and the exact finite witnesses.
The proof uses the complete applied Stickelberger and Kummer
source clauses, with their foundational class-field inputs retained.
It requires no completed full class-number computation or
GRH-conditional candidate for any of the degree-twenty-four fields.
No vanishing of an unused class eigenspace is inferred.
The remaining mixed arithmetic and $\mathcal R_{\rm cub}$/H04.4
conditions are not closed by this theorem.
